\documentclass[11pt]{article}
\usepackage[margin=1in]{geometry}
\usepackage{amsmath,amssymb,amsthm,mathtools}
\usepackage{hyperref}
\hypersetup{colorlinks,linkcolor=blue!60!black,citecolor=blue!60!black,urlcolor=blue!60!black}
\usepackage[T1]{fontenc}
\usepackage{microtype}
\usepackage{enumitem}
\setlist{nosep}

\newtheorem{remark}{Remark}
\newtheorem{proposition}{Proposition}

\newcommand{\N}{\mathbb{N}}
\newcommand{\R}{\mathbb{R}}
\newcommand{\E}{\mathbb{E}}
\newcommand{\PP}{\mathbb{P}}
\newcommand{\supnorm}{M}
\newcommand{\eps}{\varepsilon}

\title{Formalization notes for Erdős Problem 524\\ \large Remarks accumulated during the Lean 4 proof of Chojecki (2026)}
\author{Companion memo to \texttt{FormalConjectures/ErdosProblems/524.lean}}
\date{\today}

\begin{document}
\maketitle

\begin{abstract}
This memo accumulates mathematical remarks, subtleties, and open concerns
encountered while formalizing Chojecki's resolution of Erdős problem 524
(``supremum of random polynomials on $[-1,1]$'') in Lean 4. It is intended as
a running appendix to the Lean file and is updated as the formalization
progresses. The memo records where the Lean statement requires additional
hypotheses that are implicit in Chojecki's paper, and where genuine gaps in
prerequisite theory (Mathlib-level) remain.
\end{abstract}

%---------------------------------------------------------------------
\section{Setup and architecture}

Let $(a_k)_{k \ge 1}$ be i.i.d.\ Rademacher signs on a probability space
$(\Omega, \PP)$ and write
\[
  P_n(\omega)(x) = \sum_{k=1}^n a_k(\omega)\, x^k, \qquad
  \supnorm_n(\omega) = \sup_{x \in [-1,1]} |P_n(\omega)(x)|.
\]
The Lean formalization is organized around six \emph{helper modules} under
\texttt{Helpers/} that provide purely combinatorial / analytic lemmas
(\texttt{Stirling}, \texttt{CentralBinom}, \texttt{BlockIndep},
\texttt{IndepSetBridge}, \texttt{LilNormAsymptotics},
\texttt{CentralBinomWindowSum}), and a single file \texttt{524.lean} that
develops the probabilistic content.

Closed sorries (as of the latest commit on \texttt{add-erdos-524}):
\begin{itemize}
  \item \texttt{\#2 lil\_tail\_lower\_at\_scale} — sharp Rademacher lower tail
    at the LIL scale;
  \item \texttt{\#3 lil\_sparse\_lower\_bc} — block-BC on exponential
    subsequences;
  \item \texttt{\#4 kolmogorov\_lil\_lower\_bound} — full $\liminf / \limsup
    \ge 1$ half of Kolmogorov's LIL;
  \item \texttt{\#1 erdos\_524} wrapper — reduces to \texttt{sharp\_upper\_envelope};
  \item \texttt{\#5 sparse\_lower\_envelope} — conditionally, via axiomatization
    of the Chojecki ingredients.
\end{itemize}

Remaining:
\begin{itemize}
  \item \texttt{chojecki\_sparse\_lower\_envelope\_proof} — the ``inner
    assembly'' (Borel–Cantelli + block independence + cubic subsequence
    asymptotics), conditionally on the atomic sub-axioms;
  \item \texttt{\#6 exact\_lower\_constant} — open in mathematics
    (Gao–Li–Wellner exact constant is an active research problem).
\end{itemize}

%---------------------------------------------------------------------
\section{The atomic axiom set}

The bundled Chojecki axiom has been decomposed into six atomic sub-axioms
(plus two ``post-KMT transfer'' axioms stated directly in terms of
$\supnorm_n$):
\begin{description}[labelwidth=2cm,leftmargin=2.2cm]
  \item[\texttt{wiener\_process\_exists}] Existence of a standard Brownian
    motion on $[0,\infty)$;
  \item[\texttt{ito\_integral\_exp\_kernel}] Existence of the Itô integral
    $Y(u) = \int_0^1 e^{-us}\,dB(s)$ as a centered Gaussian process;
  \item[\texttt{gao\_li\_wellner\_small\_ball\_upper}] GLW upper bound
    $\PP(\sup_u |Y(u)| \le \eps) \le \exp(-\underline{c}|\log\eps|^3)$;
  \item[\texttt{gao\_li\_wellner\_small\_ball\_lower}] GLW lower bound
    $\PP(\sup_u |Y(u)| \le \eps) \ge \exp(-\bar c|\log\eps|^3)$;
  \item[\texttt{two\_dim\_KMT\_coupling}] 2D Komlós–Major–Tusnády strong
    invariance for the reparametrized pair;
  \item[\texttt{endpoint\_reparametrization}] The deterministic substitution
    $x = \pm e^{-u/n}$ that maps $\supnorm_n / \sqrt n$ to a supremum of two
    processes $Z_n^\pm$ on $u \ge 0$.
\end{description}
Transferred to the polynomial level:
\begin{description}[labelwidth=2cm,leftmargin=2.2cm]
  \item[\texttt{polynomial\_sup\_small\_ball\_upper}]
    $\PP(\supnorm_n \le \eps\sqrt n) \le \exp(-\bar c|\log\eps|^3)$;
  \item[\texttt{polynomial\_sup\_small\_ball\_lower}]
    $\exp(-\underline{c}|\log\eps|^3) \le \PP(\supnorm_n \le \eps\sqrt n)$.
\end{description}

%---------------------------------------------------------------------
\section{Remarks and subtleties}

\begin{remark}[\textbf{Naming convention for the GLW constants}]
The \texttt{GaoLiWellnerConstants} structure declares fields
\texttt{lower, upper} with docstring
``\texttt{exp(-$\bar c$|log $\eps$|$^3$) $\le$ $\PP$ $\le$ exp(-$\underline{c}$|log $\eps$|$^3$)}''
i.e.\ \texttt{upper} $= \bar c$ (the \emph{larger} constant) is associated
with the \emph{lower} bound on the probability, and \texttt{lower} $= \underline{c}$
(the smaller constant) is associated with the \emph{upper} bound.

However, the current transferred axioms
\texttt{polynomial\_sup\_small\_ball\_\{upper,lower\}} use the fields
\emph{swapped}:
\[
  \PP \le \exp(-\texttt{glw.upper}\cdot\ldots), \qquad
  \exp(-\texttt{glw.lower}\cdot\ldots) \le \PP.
\]
This is inconsistent with the structure's docstring convention and would
force $\texttt{glw.upper} \le \texttt{glw.lower}$ once both are used, contradicting
\texttt{lower\_le\_upper}. For the trivial instance
$(\texttt{lower},\texttt{upper}) = (1,1)$ the axioms are mutually consistent,
but for any strict instance the current axioms assert a claim stronger than
Gao–Li–Wellner actually provides. \emph{Recommendation:} swap the fields in
the axioms (or swap the structure's docstring convention) before attempting
to close the downstream BC argument.
\end{remark}

\begin{remark}[\textbf{The factor $6 = 3 \cdot 2$ in $\alpha_\pm$}]
Chojecki's Theorem 18 claims, along the subsequence $n_m = \lfloor e^{m^3}
\rfloor$,
\[
  \alpha_- \le \limsup_{m\to\infty}
    \frac{\log\big(\sqrt{n_m}/\supnorm_{n_m}(\omega)\big)}{(\log\log n_m)^{1/3}}
    \le \alpha_+ \quad \text{a.s.},
\]
with $\alpha_\pm = (1/(6\, c_{\mp}))^{1/3}$. The factor $6$ decomposes as $3
\cdot 2$:
\begin{itemize}
  \item the $3$ comes from the cubic subsequence asymptotic
    $\log\log n_m \sim 3\log m$ (so the BC1 threshold scales as
    $m^{-3c\beta^3}$);
  \item the $2$ comes from the endpoint split
    $\supnorm_n = \max(\supnorm_n^+, \supnorm_n^-)$ and the
    \emph{asymptotic independence} of the two halves once coupled (via
    KMT) to two independent copies of $Y$.
\end{itemize}
The second factor is a substantive input: it is \emph{not} a direct
consequence of the one-sided \texttt{polynomial\_sup\_small\_ball\_upper}
axiom. It requires an additional statement — essentially
\[
  \PP(\supnorm_n \le \eps\sqrt n) \le \PP(\supnorm_n^+ \le \eps\sqrt n)\cdot
    \PP(\supnorm_n^- \le \eps\sqrt n)(1 + o(1)),
\]
a ``product-of-marginals'' small-ball bound for the two halves. Without
encoding this separately (or a sharper GLW asymptotic that already
accounts for the factor), the BC1 direction only closes under the
structural hypothesis
\begin{equation}\label{eq:gap}
  2\,\texttt{glw.lower} \le \texttt{glw.upper},
\end{equation}
which is currently added as a new field \texttt{two\_lower\_le\_upper} to
\texttt{GaoLiWellnerConstants}. This gap condition is sufficient for the
present axiomatic setup; it is implied by (but weaker than) the true
factor-of-$2$ independence argument.
\end{remark}

\begin{remark}[\textbf{Block independence and $\supnorm_n$}]
The helper \texttt{iIndepFun\_block\_sums} gives mutual independence of
\emph{block sums} $Y_k = \sum_{j \in I_k} a_j$ over pairwise disjoint index
blocks $I_k$. This is directly applicable to the LIL lower bound
(\texttt{\#4}) because the walk increment $S_{n_{k+1}} - S_{n_k}$ \emph{is}
such a block sum.

However, the polynomial supremum $\supnorm_{n_m}(\omega)$ depends on the
full prefix $(a_1,\ldots,a_{n_m})$, not on an increment. Consequently, the
events $\{\supnorm_{n_m} \le \eps_m\sqrt{n_m}\}$ are \emph{not} mutually
independent via block-independence alone. Chojecki's proof resolves this by
(a) using KMT to couple $\supnorm_n / \sqrt n$ to a functional of the
Brownian motion $B$ on $[0,1]$, then (b) decomposing $B$ into independent
increments over the cubic subsequence and truncating outside a block. The
truncation argument is \emph{not} a pure block-independence instance; it
is a coupling + truncation argument which is currently bundled inside the
\texttt{chojecki\_sparse\_lower\_envelope\_proof} sorry.

Closing the proof rigorously in Lean requires either an additional axiom
of the form ``block-truncated $\supnorm_{n_m}^\star$ are independent and
$\ell^\infty$-close to $\supnorm_{n_m}$ on a set of full measure'', or a
direct formalization of the KMT + truncation argument. The latter is a
substantial multi-month effort.
\end{remark}

\begin{remark}[\textbf{Sharpness of the LIL formalization (sorry \#2)}]
The sharp Rademacher lower tail
\[
  \PP(S_n \ge (1-\delta)\sqrt{2n\log\log n})
    \ge C\cdot \exp\big(-((1-\delta)^2 + \delta)\log\log n\big)
\]
is achieved with the polynomial Stirling bound refined to second-order
Kullback–Leibler. With the naive first-order bound $\log x \le x-1$, the
central binomial PMF yields only Gaussian constant $\alpha = 4$, which is
\emph{provably insufficient}: for any $\delta < (3-\sqrt5)/2 \approx 0.382$,
the implied bound $\propto (\log n)^{-2(1-\delta)^2}$ fails to dominate
the target $(\log n)^{-((1-\delta)^2+\delta)}$.

Fix: sharpen the KL bound via a second-order Taylor expansion of
$\log(1\pm y)$, restricted to a regime $|k - n/2|^2 \le n\log n$. This
regime constraint is automatically satisfied at the LIL scale
$|k-n/2| \le \sqrt{n\log\log n}$. With the sharp bound,
$\alpha = 2+\eps$ is obtainable with $\eps = \delta/4$, and the target
exponent $(1-\delta)^2 + \delta$ is comfortably met.

\emph{Tactical note.} The algebraic core \texttt{margin\_absorbs\_sqrt\_core}
introduces $u = \sqrt{2s}$ as a hypothesis $u^2 = 2s$, keeping all proof
bodies polynomial in $(u, s, \delta)$. This sidesteps the heartbeat
blowup that defeated three earlier attempts using $\texttt{Real.sqrt}$
directly.
\end{remark}

\begin{remark}[\textbf{On \texttt{answer(sorry)}}]
The FormalConjectures framework elaborates
\texttt{answer(sorry)} in a \texttt{Prop} context to \texttt{True} by default
(see \texttt{FormalConjectures/Util/Answer.lean}, option
\texttt{google.answer = alwaysTrue}). Consequently the theorem
\texttt{erdos\_524 : answer(sorry) $\leftrightarrow$ \emph{RHS}} is
semantically $\texttt{True} \leftrightarrow \emph{RHS}$, i.e., the RHS is
the actual mathematical content. The upper envelope order of magnitude
$M_n(\omega) \asymp \sqrt{n\log\log n}$ a.s.\ is proved (and therefore the
wrapper is closable) via the combined limsup result
$\limsup \supnorm_n / \sqrt{2n\log\log n} = 1$ a.s., which is now
sorry-free in the repo.
\end{remark}

%---------------------------------------------------------------------
\section{Change log}

\begin{itemize}
  \item \textbf{2026-04-23, commit \texttt{6e7bc7e}.} Refactored the big
    \texttt{chojecki\_sparse\_lower\_envelope} axiom into six atomic
    sub-axioms plus two transferred axioms. Closed sorries \#1 and \#5
    (the latter conditionally, via the new sub-axioms).
  \item \textbf{2026-04-23, commit \texttt{315e8e1}.} Closed sorries \#2,
    \#3, \#4. Added six helper modules (\texttt{StirlingTwoSided},
    \texttt{IndepSetBridge}, \texttt{BlockIndep},
    \texttt{LilNormAsymptotics}, \texttt{CentralBinomLower},
    \texttt{CentralBinomWindowSum}). Added \texttt{walk\_pmf\_binomial} in
    524.lean.
  \item \textbf{2026-04-23 (uncommitted).} Added field
    \texttt{two\_lower\_le\_upper : 2·lower $\le$ upper} to
    \texttt{GaoLiWellnerConstants}, updating the default instance to
    $(\texttt{lower},\texttt{upper}) = (1,2)$. This unblocks the BC1 upper
    half of \texttt{chojecki\_sparse\_lower\_envelope\_proof}.
\end{itemize}

\end{document}
